\documentclass[11pt]{article}
\usepackage[T1]{fontenc}
\usepackage{lmodern}
\usepackage{amsmath,amssymb,amsthm}
\usepackage[margin=1in]{geometry}
\usepackage{microtype}
\usepackage{hyperref}
\hypersetup{colorlinks=true,linkcolor=blue,citecolor=blue,urlcolor=blue,
  pdftitle={The location of the most isolated point in a polyhedral sample},
  pdfauthor={Kaiwen Jin}}
\newtheorem{theorem}{Theorem}
\newtheorem{lemma}[theorem]{Lemma}
\newtheorem{proposition}[theorem]{Proposition}
\newcommand{\PP}{\mathbb P}
\newcommand{\EE}{\mathbb E}
\newcommand{\HH}{\mathcal H}
\newcommand{\norm}[1]{\lVert #1\rVert}
\DeclareMathOperator{\dist}{dist}
\title{The location of the most isolated point\ in a polyhedral sample}
\author{Kaiwen Jin\\\small\href{mailto:yui.ui.siki@gmail.com}{yui.ui.siki@gmail.com}}
\date{October 14, 2026}

\begin{document}
\maketitle
\begin{abstract}
Let $A$ be a compact convex polyhedron in $\mathbb R^3$ with nonempty
interior, and sample independent uniform points in $A$. We prove that the
point with the largest nearest-neighbour distance converges in distribution
to normalized surface area on $\partial A$ when every interior dihedral
angle exceeds $\pi/2$. When the smallest angle is less than $\pi/2$, its
limit is normalized length on the union of all edges with that angle.
The proof computes spatial empty-ball intensities and joint exceedance
probabilities for the fixed-size sample. The critical angle is excluded.
\end{abstract}

\section{The location question}

In a geometric graph with connection radius $r$, a sample point is isolated
precisely when its nearest-neighbour distance exceeds $r$. The point whose
isolation persists longest therefore has the largest such distance. The
boundary affects both the size of this distance and the location of the
point attaining it: an empty ball costs less near a face or an edge, but
fewer sample points are available there.

Let $A\subset\mathbb R^3$ be a fixed compact convex polyhedron with nonempty
interior, and write $V=|A|$ for its volume. Let $X_1,X_2,\ldots$ be
independent uniform points in $A$. For $n\geq2$ put
\[
 R_{n,i}=\min_{\substack{1\leq j\leq n\\j\ne i}}\norm{X_i-X_j},
 \qquad L_n=\max_{1\leq i\leq n}R_{n,i},
 \qquad I_n=\min\{i:R_{n,i}=L_n\},\qquad Z_n=X_{I_n}.
\]
For an edge $e$, let $\alpha_e\in(0,\pi)$ denote its interior dihedral
angle, and set
\[
 \alpha=\min_e\alpha_e,\qquad
 E_* =\bigcup_{\alpha_e=\alpha}e,\qquad
 S=\HH^2(\partial A),\qquad L=\HH^1(E_*).
\]
Here $\HH^j$ has the usual Euclidean normalization, so these are surface
area and edge length. Both $S$ and $L$ are positive and finite.

Penrose, Yang and Higgs proved the leading strong law
\[
 \frac{nL_n^3}{V\log n}\longrightarrow
 \max\left\{\frac1\pi,\frac1{2\alpha}\right\}
 \quad\hbox{almost surely}
\]
and proposed the corresponding spatial limits
\cite[Introduction, p.~1724; Corollary~2.3]{PYH}.
Their strong law identifies the competition between faces and sharp edges,
but does not determine the spatial distribution of the maximizer. Later
fluctuation results of Penrose and Yang treat smooth domains and planar
polygons \cite[arXiv v3, Theorem~1.1]{PY}; general three-dimensional polyhedra are
outside that theorem's hypotheses. We prove the two proposed spatial limits.

\begin{theorem}\label{thm:location}
For every compact convex polyhedron $A\subset\mathbb R^3$ with nonempty
interior, the point $Z_n$ defined above satisfies
\[
 Z_n\ \xrightarrow{\ d\ }\
 \begin{cases}
   \HH^2|_{\partial A}/S,&\alpha>\pi/2,\\[2mm]
   \HH^1|_{E_*}/L,&\alpha<\pi/2.
 \end{cases}
\]
The statement includes all minimum-angle edges and the specified
smallest-index rule for ties. The case $\alpha=\pi/2$ is not considered.
\end{theorem}

We first compute the intensity of points with an empty neighbourhood at an
extreme radius. Its spatial weight is constant per unit area on a facet,
and constant per unit length on edges with the same angle. To identify the
maximizer, we also need joint probabilities. A convex-geometric volume
increment makes nearby simultaneous extreme events negligible; factorial
moments then give joint no-exceedance limits. Comparing maxima in a set and
its complement proves Theorem~\ref{thm:location}.

All estimates below concern this fixed $A$. Constants may also depend on a
fixed finite set of thresholds or a fixed moment order. Balls are Euclidean
and closed. We use the elementary product-integration and convergence
results in \cite[Sections~1.5, 1.7 and 3.10]{Durrett}, and the linear
volume transformation formula in \cite[Section~2.8]{Hunter}.

\section{Local polyhedral geometry}

For $x\in A$ and $r>0$, write
\[
 v_r(x)=|A\cap B(x,r)|.
\]
Two estimates will control the vertices and the dependence between nearby
empty-ball events.

\begin{lemma}\label{lem:union}
There is $b>0$ such that $v_r(x)\geq br^3$ for every $x\in A$ and all
sufficiently small $r$. Moreover, for every fixed $M\geq1$ there are
$c>0$ and $r_0>0$ such that
\begin{equation}\label{eq:union}
 |A\cap(B(x,r)\cup B(y,r))|
 \geq\min\{v_r(x),v_r(y)\}+cr^3
\end{equation}
whenever $0<r<r_0$ and $r\leq\norm{x-y}\leq Mr$.
\end{lemma}

\begin{proof}
Choose $z\in A$, $\rho>0$ and $D\geq\sup_{x\in A}\norm{z-x}+\rho$ with
$B(z,\rho)\subset A$. Convexity gives
\begin{equation}\label{eq:inball}
 B(x+u(z-x),u\rho)\subset A\qquad(0\leq u\leq1).
\end{equation}
For $u=r/(2D)$ this ball is contained in $B(x,r)$, proving the first
assertion with $b=(4\pi/3)(\rho/(2D))^3$.

For the second assertion let $d=\norm{x-y}$, $q=(x+y)/2$ and
$u=(y-x)/d$. If $\norm{q-z}<\rho/2$, then for $r$ sufficiently small
both radius-$r$ balls lie inside $B(z,\rho)$. The ball
$B(y+(r/2)u,r/8)$ lies in $B(y,r)$ and outside $B(x,r)$.

Suppose instead that $\norm{q-z}\geq\rho/2$. Choose $w=y$ if
$u\cdot(z-q)\geq0$, and $w=x$ otherwise. Let $w'$ be the other point,
$u_w=(w-w')/d$, $\theta=r/(4D)$, and $p=w+\theta(z-w)$. Then
$u_w\cdot(z-q)\geq0$, $B(p,\theta\rho)\subset A$, and $\norm{p-w}\leq r/4$.
Also
\begin{align*}
 \norm{p-w'}^2
 &=\norm{(1-\theta/2)d u_w+\theta(z-q)}^2\\
 &\geq(1-\theta)d^2+\theta^2\rho^2/4\\
 &\geq r^2-\frac{M^2r^3}{4D}+\frac{\rho^2r^2}{64D^2}.
\end{align*}
For small $r$, this is at least $(1+\rho^2/(128D^2))r^2$.
Hence $\norm{p-w'}\geq r+\eta_0r$ for some $\eta_0>0$ depending on $A$.
Choose
$0<\eta<\min\{\eta_0/2,\rho/(8D),1/8\}$.
The ball $B(p,\eta r)$ lies in $A\cap B(w,r)$ and outside $B(w',r)$.

In either case the union contains one of the clipped balls and an
exclusive ball of volume at least a fixed positive multiple of $r^3$.
Subtracting the smaller individual volume proves \eqref{eq:union}.
\end{proof}

We will also use \eqref{eq:union} for radii between $r_-$ and $r_+$ with
$r_+/r_-\to1$. The volume change in increasing a radius from $r_-$ to
$r_+$ is at most $(4\pi/3)(r_+^3-r_-^3)=o(r_-^3)$. Consequently, if
$\max\{r_i,r_j\}\leq\norm{x-y}\leq Mr_+$, then for large $n$
\begin{equation}\label{eq:mixed-union}
 |A\cap(B(x,r_i)\cup B(y,r_j))|
 \geq\min\{v_{r_i}(x),v_{r_j}(y)\}+c'r_-^3.
\end{equation}
Here Lemma~\ref{lem:union} is applied at $r_-$, with a fixed enlargement
of $M$, and half its increment absorbs the shell-volume error.

We next partition $A$ into regions where a small ball sees a full space,
a half-space or a wedge. Let $\mathcal F$, $\mathcal E$ and $\mathcal V$
be the facets, edges and vertices of $A$. For $f\in\mathcal F$, let
$d_f$ be the affine inward signed distance to its supporting plane.
Thus $A=\{x:d_f(x)\geq0\text{ for all }f\in\mathcal F\}$, with
irredundant facet inequalities.

\begin{lemma}\label{lem:layers}
There are $C>0$ and $r_0>0$ such that, for $0<r<r_0$, outside
$\bigcup_{v\in\mathcal V}B(v,Cr)$ the set
\[
 J_r(x)=\{f\in\mathcal F:d_f(x)<r\}
\]
is empty, a singleton, or the two facets incident to an edge.
The singleton region for $f$ lies between the coordinate domains
$f'_r\times[0,r]$ and $f\times[0,r]$, with
$\HH^2(f\setminus f'_r)=O(r)$. The two-facet region for $e$ lies between
$e'_r\times[0,r]^2$ and $e\times[0,r]^2$, with
$\HH^1(e\setminus e'_r)=O(r)$. In the edge coordinates $(y,h_1,h_2)$,
where $y\in e$ is orthogonal projection and $h_i$ is inward face distance,
\begin{equation}\label{eq:jacobian}
 dx=\frac{d\HH^1(y)\,dh_1\,dh_2}{\sin\alpha_e}.
\end{equation}
All domain inclusions are understood up to null boundaries.
\end{lemma}

\begin{proof}
A point of a positive-dimensional face is a convex combination of its
vertices: a maximal line segment through a point in its minimal face has
endpoints on proper subfaces, and induction on face dimension gives the
assertion. The vertices are finite because the facet planes are finite.

For a nonempty subset $J\subset\mathcal F$, let
$F_J=A\cap\{d_f=0:f\in J\}$ and $D_J=\sum_{f\in J}d_f$.
If $F_J\ne\varnothing$, $D_J$ vanishes on its vertices and is strictly
positive at every other vertex. Let $\delta_J>0$ be the minimum of these
positive values. In a vertex convex combination for $x$, the total weight
$\beta$ on vertices outside $F_J$ is at most $D_J(x)/\delta_J$.
When this is less than one, normalizing the remaining weights gives a
point of $F_J$ within distance $\operatorname{diam}(A)\beta$ of $x$.
Thus
\[
 \dist(x,F_J)\leq C_JD_J(x)\qquad\text{when }D_J(x)\text{ is small}.
\]
If $F_J$ is empty, $D_J$ has a positive minimum on compact $A$.
There are finitely many $J$, so an active subset $J_r(x)$ has nonempty
common face and puts $x$ within $O(r)$ of that face.

Outside sufficiently large $Cr$ vertex balls, the common face has
positive dimension. A facet is specified by a single irredundant
inequality. An edge has exactly two incident maximal facets, since a
transverse slice at a relative-interior edge point is a planar wedge
with two boundary rays. Three distinct maximal facets cannot contain an
edge. This proves the asserted alternatives and includes nonsimple vertices.

For a singleton $f$, write $x=y+h u_f$ with $u_f$ its inward unit normal
and $0\leq h<r$. Every other slack is at least $r$, so
$d_g(y)\geq d_g(x)-h\geq0$ and $y\in f$.
For a lower domain take those $y\in f$ with $d_g(y)\geq2r$ for all
$g\ne f$ and $\dist(y,\mathcal V)\geq(C+1)r$.
These conditions ensure that only $f$ is active and the vertex balls
are avoided. Each nonconstant affine restriction $d_g|_f$ removes a
strip of width $O(r)$ from a bounded polygon; constant restrictions are
strictly positive. Vertex disks remove $O(r^2)$ area. The total area lost is
$O(r)$.

For an edge $e$, the map from transverse orthonormal coordinates to the
two inward face distances has determinant of absolute value
$\sin\alpha_e$, because the inward normals have angle $\pi-\alpha_e$.
This proves \eqref{eq:jacobian}. If $0\leq h_i<r$, the transverse
displacement is at most $C_er$. The face-distance bound also gives
$\dist(x,e)\leq C'_er$. A projection beyond an endpoint would therefore
put $x$ in an $O(r)$ vertex ball; enlarge $C$ to exclude it.
For a lower domain trim $O(r)$ length from the ends of $e$, so all
nonincident slacks at $y$ exceed $(C_e+1)r$ and
$\dist(y,\mathcal V)>(C+C_e)r$. Each such slack is affine and positive
in the relative interior, so the required trimming has size $O(r)$.
The transverse displacement preserves slack at least $r$ and distance
greater than $Cr$ from the vertices. This proves the lower inclusion.
All constants may be chosen simultaneously over the finitely many faces.
\end{proof}

\section{Spatial empty-ball intensities}

Put $\lambda_n=n/V$. The conditional probability that a sample point
at $x$ has no other point within distance $r$ is
$(1-v_r(x)/V)^{n-1}$. At the radii used below it is uniformly asymptotic
to $e^{-\lambda_nv_r(x)}$. We begin with the latter expression.

For fixed $s\in\mathbb R$, define $r_n(s)>0$ for sufficiently large $n$
by $t_n(s)=\lambda_nr_n(s)^3$, where
\begin{equation}\label{eq:thresholds}
 t_n(s)=
 \begin{cases}
  \pi^{-1}(\log n-\log\log n)+(3/(2\pi))s,&\alpha>\pi/2,\\[1mm]
  (2\alpha)^{-1}(\log n-4\log\log n)+(3/(2\alpha))s,&\alpha<\pi/2.
 \end{cases}
\end{equation}
Set
\begin{equation}\label{eq:weights}
 (\nu,K)=
 \begin{cases}
  \left(\HH^2|_{\partial A}/S,\ S/(\pi^{1/3}V^{2/3})\right),&\alpha>\pi/2,\\[1mm]
  \left(\HH^1|_{E_*}/L,\ 4(2\alpha)^{4/3}L/(\pi^2\sin\alpha\,V^{1/3})\right),&\alpha<\pi/2.
 \end{cases}
\end{equation}

\begin{proposition}\label{prop:intensity}
For every fixed $s\in\mathbb R$, the finite measures on $A$
\[
 \mu_{n,s}(dx)=\lambda_ne^{-\lambda_nv_{r_n(s)}(x)}\,dx
\]
converge weakly to $Ke^{-s}\nu$. The same limit holds on every
$\nu$-continuity Borel set $B\subset A$ for the expected number of indices
with $X_i\in B$ and $R_{n,i}>r_n(s)$.
\end{proposition}

\begin{proof}
Let $r=r_n(s)$, $t=\lambda_nr^3$, and let $\psi$ be continuous on $A$.
In a half-space the unit-ball volume at inward center distance $u\in[0,1]$
is
\[
 g_F(u)=\frac{2\pi}3+\pi u-\frac{\pi}3u^3.
\]
This follows by integrating the disk sections from height $-u$ to $1$.
Since $g_F(u)-g_F(0)\geq(2\pi/3)u$ and $g_F'(0)=\pi$, rescaling
$u=w/t$ gives
\begin{equation}\label{eq:face-laplace}
 \int_0^1e^{-tg_F(u)}\,du
 \sim\frac{e^{-2\pi t/3}}{\pi t}.
\end{equation}
The rescaled integrand is dominated on $\mathbb R_+$ by
$e^{-(2\pi/3)w}$ and tends to $e^{-\pi w}$.

For an edge of angle $\beta\in(0,\pi)$ let $a_1,a_2$ be its inward
unit normals and define
\[
 g_\beta(u_1,u_2)=
 \left|\{z\in B(0,1):a_i\cdot z\geq-u_i,\ i=1,2\}\right|.
\]
At $(0,0)$ the wedge occupies fraction $\beta/(2\pi)$ of the ball, so
$g_\beta(0,0)=2\beta/3$. Relaxing one inequality adds a slab whose
cross-sectional area tends to half a unit disk, $\pi/2$.
The remaining doubly negative region has third-coordinate length at most
$2$ and face-distance area at most $u_1u_2/\sin\beta$. It follows that
\begin{equation}\label{eq:wedge-expansion}
 g_\beta(u_1,u_2)=\frac{2\beta}3+
 \frac\pi2(u_1+u_2)+o(u_1+u_2).
\end{equation}
Slice-area continuity follows from dominated convergence, since the limiting
straight-line boundary has area zero. The function $g_\beta$ is continuous
and coordinatewise nondecreasing, and an increase of either coordinate
from zero adds positive volume. Compactness of $[0,1]^2$ and
\eqref{eq:wedge-expansion} give
\[
 g_\beta(u_1,u_2)-\frac{2\beta}3\geq c_\beta(u_1+u_2)
 \qquad(0\leq u_i\leq1)
\]
for some $c_\beta>0$. Rescaling both coordinates by $t$ and using this
integrable exponential bound yields
\begin{equation}\label{eq:edge-laplace}
 \int_{[0,1]^2}e^{-tg_\beta(u_1,u_2)}\,du_1du_2
 \sim\frac{4e^{-2\beta t/3}}{\pi^2t^2}.
\end{equation}

Apply Lemma~\ref{lem:layers}. Replacing $\psi(x)$ by its facet or edge
projection changes it by at most its modulus of continuity at $C_Ar$.
The upper and lower domains differ by $O(r)$ area or length, multiplied
by the same positive transverse integral. Thus the contribution of a
facet $f$ is
\begin{equation}\label{eq:face-contribution}
 \frac{\lambda_nr e^{-2\pi t/3}}{\pi t}
 \left(\int_f\psi\,d\HH^2+o(1)\right)(1+o(1)),
\end{equation}
and that of an edge $e$ of angle $\beta$ is
\begin{equation}\label{eq:edge-contribution}
 \frac{4\lambda_nr^2 e^{-2\beta t/3}}{\pi^2\sin\beta\,t^2}
 \left(\int_e\psi\,d\HH^1+o(1)\right)(1+o(1)).
\end{equation}
These estimates also hold for signed $\psi$ by the same absolute error
bound. The full-ball region contributes at most
$\norm{\psi}_\infty n e^{-4\pi t/3}$ in absolute value. The vertex
region has volume $O(r^3)$, and Lemma~\ref{lem:union} bounds its contribution
by $C_A\norm{\psi}_\infty t e^{-bt}$.

In the face regime, $t\sim\pi^{-1}\log n$. The prefactor in
\eqref{eq:face-contribution} tends to $\pi^{-1/3}V^{-2/3}e^{-s}$.
The full-ball term is $O_s(n^{-1/3}(\log n)^{4/3})$, and the vertex term
tends to zero. For an edge of angle $\beta$, its prefactor in
\eqref{eq:edge-contribution} is a constant times
\[
 n^{1/3-2\beta/(3\pi)}(\log n)^{-4/3+2\beta/(3\pi)}.
\]
Every $\beta\geq\alpha>\pi/2$, so these terms vanish. Summing the facet
integrals gives $Ke^{-s}\int\psi\,d\nu$ with the first row of
\eqref{eq:weights}.

In the edge regime, $t\sim(2\alpha)^{-1}\log n$. A minimum-angle edge
has prefactor tending to
$4(2\alpha)^{4/3}e^{-s}/(\pi^2\sin\alpha\,V^{1/3})$.
For another edge the power of $n$ is $1/3-\beta/(3\alpha)<0$.
For facets it is $2/3-\pi/(3\alpha)<0$, and for the full-ball region it
is $1-2\pi/(3\alpha)<0$. All accompanying logarithmic powers are fixed.
The vertex term again vanishes. Only edges of angle $\alpha$ remain,
and their identical prefactors give the second row of \eqref{eq:weights}.
This proves the finite-measure weak limit.

Finally, $v_r(x)\leq(4\pi/3)r^3=O(\log n/n)$ uniformly. For fixed $k$,
\begin{equation}\label{eq:binomial-approximation}
 (1-v_r(x)/V)^{n-k}
 =e^{-\lambda_nv_r(x)}(1+o(1)),
\end{equation}
uniformly, since the logarithm error is $O_k((\log n)^2/n)$.
Taking $k=1$ identifies the actual first moments. The total mass of
$\mu_{n,s}$ tends to $Ke^{-s}>0$; normalizing it gives probability
measures converging weakly to $\nu$. The continuity-set characterization
of weak convergence, followed by multiplication by the total mass, proves
the asserted limits on $B$.
\end{proof}

\section{Joint no-exceedance limits}

We now keep the spatial categories while calculating the joint events in
the original sample of exactly $n$ points.

\begin{proposition}\label{prop:void}
Let $B_1,\ldots,B_m$ be pairwise disjoint $\nu$-continuity Borel subsets of
$A$, and let $s_1,\ldots,s_m\in\mathbb R$ be fixed. If $W_n$ counts the
indices $i$ for which $X_i\in B_j$ and $R_{n,i}>r_n(s_j)$ for some $j$, then
\begin{equation}\label{eq:void}
 \PP(W_n=0)\longrightarrow
 \exp\left(-K\sum_{j=1}^m\nu(B_j)e^{-s_j}\right).
\end{equation}
\end{proposition}

\begin{proof}
Fix $k\geq1$ and take $n>k$. Expand the falling factorial moment $\EE[(W_n)_k]$ over
the spatial categories of $k$ ordered selected indices. For one such
assignment let $r_i$ be the radius of the category of $x_i$. Conditional
on these locations, the joint exceedance event is impossible unless
\begin{equation}\label{eq:compatibility}
 \norm{x_i-x_l}>\max\{r_i,r_l\}\qquad(i\ne l).
\end{equation}
When this holds, the remaining $n-k$ points must avoid the ball union,
so the corresponding factorial integral is
\begin{equation}\label{eq:factorial-integral}
 \frac{(n)_k}{V^k}\int_{B_{j_1}\times\cdots\times B_{j_k}}
 \mathbf1_{\eqref{eq:compatibility}}
 \left(1-\frac{|A\cap\bigcup_{i=1}^kB(x_i,r_i)|}{V}\right)^{n-k}
 dx_1\cdots dx_k.
\end{equation}
Here $(w)_k=w(w-1)\cdots(w-k+1)$ and $(w)_0=1$ for nonnegative integers $w$. The indicator in \eqref{eq:factorial-integral} enforces \eqref{eq:compatibility}. Let $r_-$ and $r_+$ be the smallest
and largest of the finitely many $r_n(s_j)$. Their ratio tends to one,
and $\lambda_nr_\pm^3$ is asymptotic to a positive multiple of $\log n$.

If all pair distances exceed $2r_+$, the balls are disjoint. Their union
volume is $O_k(\log n/n)$, so the same logarithm estimate as in
\eqref{eq:binomial-approximation} changes the probability in
\eqref{eq:factorial-integral}, uniformly, to
\[
 (1+o(1))\prod_{i=1}^k e^{-\lambda_nv_{r_i}(x_i)}.
\]
The unrestricted product integral has limit
$\prod_i Ke^{-s_{j_i}}\nu(B_{j_i})$, by
Proposition~\ref{prop:intensity} and $(n)_k/n^k\to1$.
The part removed from this product by a pair at distance at most $2r_+$
is $o(1)$. Indeed one point has a scaled integral $O(1)$, its partner
is integrated over volume $O(r_+^3)$ and has exponential weight at most
$e^{-b\lambda_nr_-^3}$, and all other scaled integrals are $O(1)$.
The resulting bound is
\begin{equation}\label{eq:product-near}
 C_k\lambda_nr_+^3e^{-b\lambda_nr_-^3}=o(1).
\end{equation}

It remains to bound compatible configurations with a close pair in the
actual integral \eqref{eq:factorial-integral}. Join two of the $k$ centers
when their distance is at most $2r_+$. Balls in distinct connected
components are disjoint. In each component choose a root with the smallest
individual volume $v_{r_i}(x_i)$; there are only finitely many possible root
choices. A component with at least two vertices has an edge satisfying
\eqref{eq:compatibility}. By \eqref{eq:mixed-union}, its union volume is
at least the root volume plus $c r_-^3$: the smaller volume of this pair
is at least the minimum over the component.

Every center of a component lies within $2(k-1)r_+$ of its root.
If it has $q$ vertices, integration of its other $q-1$ coordinates costs
at most $C_kr_+^{3(q-1)}$. Put $\lambda'_n=(n-k)/V$, and use
$(1-u)^{n-k}\leq e^{-(n-k)u}$ in the original conditional probability.
Each root satisfies
\[
 \lambda_n\int_Ae^{-\lambda'_nv_{r_i}(x)}\,dx=O(1),
\]
because replacing $\lambda_n$ by $\lambda'_n$ changes its integrand by
$1+o(1)$ uniformly. If there are $\ell$ components and at least one is
nontrivial, the scaled contribution is therefore at most
\begin{equation}\label{eq:cluster-bound}
 C_k(\lambda_nr_+^3)^{k-\ell}e^{-c\lambda'_nr_-^3}=o(1).
\end{equation}
For fixed $k$, the number of component partitions and root choices is
finite. The exponential in \eqref{eq:cluster-bound} is a negative power of
$n$, while its prefactor is a fixed power of $\log n$. Category and
separation constraints can be dropped when bounding the nonroot volumes
and root integrals. Summing all assignments gives, with
$a=K\sum_j\nu(B_j)e^{-s_j}$,
\[
 \EE[(W_n)_k]\longrightarrow a^k\qquad\text{for every fixed }k\geq0.
\]

For an integer $w\geq1$, Pascal's identity gives
\[
 \sum_{k=0}^h(-1)^k\binom wk=(-1)^h\binom{w-1}h,
\]
with terms above $w$ taken as zero. For $w=0$ the sum is one.
Its odd truncations bound $\mathbf1_{\{w=0\}}$ from below and its even
truncations from above. Hence, for every fixed $h\geq0$,
\[
 \sum_{k=0}^{2h+1}\frac{(-1)^k\EE[(W_n)_k]}{k!}
 \leq\PP(W_n=0)\leq
 \sum_{k=0}^{2h}\frac{(-1)^k\EE[(W_n)_k]}{k!}.
\]
First send $n$ to infinity, and then $h$ to infinity. Both exponential
series converge to $e^{-a}$, proving \eqref{eq:void}.
\end{proof}

\section{The maximizing point}

\begin{proof}[Proof of Theorem~\ref{thm:location}]
The increasing score whose level $s$ corresponds to $r_n(s)$ is
\[
 T_n(r)=
 \begin{cases}
  (2\pi/3)\lambda_nr^3-(2/3)\log n+(2/3)\log\log n,&\alpha>\pi/2,\\
  (2\alpha/3)\lambda_nr^3-(1/3)\log n+(4/3)\log\log n,&\alpha<\pi/2.
 \end{cases}
\]
Let $B$ be a $\nu$-continuity Borel subset of $A$, and put $q=\nu(B)$.
Suppose first that $0<q<1$. Let $M_{n,B}$ be the maximum of
$T_n(R_{n,i})$ over points in $B$, with empty maximum $-\infty$, and
define $M_{n,A\setminus B}$ in the same way. Proposition~\ref{prop:void}
at thresholds $s,t$ gives
\begin{equation}\label{eq:joint-maxima}
 \PP(M_{n,B}\leq s,\ M_{n,A\setminus B}\leq t)
 \longrightarrow e^{-Kq e^{-s}}e^{-K(1-q)e^{-t}}.
\end{equation}
An empty maximum in $B$ is at most $-H$ for every fixed $H>0$; its
limsup probability is consequently at most $e^{-Kq e^H}$, which tends
to zero as $H\to\infty$. The same holds for the complement. Replacing
empty maxima by zero changes any joint CDF or comparison probability by
$o(1)$. Thus the multivariate CDF criterion applies to real-valued maxima
and yields a pair of independent continuous limiting variables $G_B,G_C$
with the CDF in \eqref{eq:joint-maxima}.

The variables $e^{-G_B}$ and $e^{-G_C}$ are independent exponentials
with rates $Kq$ and $K(1-q)$. Therefore
\[
 \PP(G_B>G_C)=
 \int_0^\infty Kq\,e^{-Kq u}e^{-K(1-q)u}\,du=q,
 \qquad \PP(G_B=G_C)=0.
\]
The strict and nonstrict comparison sets in $\mathbb R^2$ have this
zero-mass diagonal boundary. Since the score is increasing,
\[
 \PP(M_{n,B}>M_{n,A\setminus B})
 \leq\PP(Z_n\in B)\leq
 \PP(M_{n,B}\geq M_{n,A\setminus B}).
\]
Both bounds tend to $q$. This retains the specified tie rule, including
the possibility that mutual nearest neighbours have the same distance.

If $q=0$, for every fixed $s$ the event $Z_n\in B$ implies either an
exceedance in $B$ above $r_n(s)$, or $L_n\leq r_n(s)$. The first probability
is bounded by its expected count, which tends to zero. The second tends to
$e^{-Ke^{-s}}$ by Proposition~\ref{prop:void} with $B_1=A$. Sending
$s\to-\infty$ after $n\to\infty$ gives $\PP(Z_n\in B)\to0$.
For $q=1$ apply this argument to the complement. Thus the probabilities
converge to $\nu(B)$ for every continuity set, proving weak convergence
on $A$ and hence on $\mathbb R^3$.
\end{proof}

\paragraph{AI-use statement.}
The mathematical argument, source checks and manuscript were prepared with
OpenAI Codex. Independent mathematical review is requested.

\begin{thebibliography}{9}
\bibitem{PYH}
M.~D. Penrose, X. Yang and F. Higgs,
\emph{Largest nearest-neighbour link and connectivity threshold in a polytopal
random sample}, Journal of Applied and Computational Topology \textbf{8}
(2024), 1723--1750.
\href{https://doi.org/10.1007/s41468-023-00154-5}{doi:10.1007/s41468-023-00154-5}.

\bibitem{PY}
M.~D. Penrose and X. Yang,
\emph{Fluctuations of the connectivity threshold and largest nearest-neighbour
link}, The Annals of Applied Probability \textbf{35}(6) (2025), 3906--3941.
\href{https://doi.org/10.1214/25-AAP2210}{doi:10.1214/25-AAP2210}.
Preprint (version 3): \url{https://arxiv.org/abs/2406.00647v3}.

\bibitem{Durrett}
R. Durrett, \emph{Probability: Theory and Examples}, fifth edition,
Cambridge University Press, 2019. Author version dated January 11, 2019:
\url{https://sites.math.duke.edu/~rtd/PTE/PTE5_011119.pdf}.

\bibitem{Hunter}
J.~K. Hunter, \emph{Measure Theory}, University of California at Davis
lecture notes, 2011.
\url{https://www.math.ucdavis.edu/~hunter/measure_theory/measure_notes_old.pdf}.
\end{thebibliography}
\end{document}
